\chapter{Appendix - Mathematical derivations}


\section{System of Equations for Known Growth Rates}
\label{app:KNWNO W}

This appendix derives the system of equations that relates observed mutation counts 
to unknown event times under the assumption that the growth rate $w$ is known and 
constant across all clones. We present the derivation for a general branch and then
illustrate the full system for a concrete example with four clones in a linear 
topology, where two copy-number alteration events occur on the branches of clones 
2 and 3 respectively.

\subsection*{General Derivation}

For a branch without copy-number alterations, the expected number of mutations 
accumulated along the branch leading to clone $i$ is
%
\begin{equation}
    \mathbb{E}[m_i] = \mu \cdot w \cdot l_{\mathrm{tot},i} \cdot (t_i - t_{\mathrm{pa}(i)}).
\end{equation}
%
Treating the observed mutation count $m_i$ as equal to its expectation yields the 
deterministic equation
%
\begin{equation}
    m_i = \mu \cdot w \cdot l_{\mathrm{tot},i} \cdot (t_i - t_{\mathrm{pa}(i)}),
\end{equation}
%
which can be solved for $t_i$ given $t_{\mathrm{pa}(i)}$:
%
\begin{equation}
    t_i = \frac{m_i}{\mu \cdot w \cdot l_{\mathrm{tot},i}} + t_{\mathrm{pa}(i)}.
\end{equation}

For a branch carrying a copy-number alteration event at time $t_{CN,ij}$, the 
mutations in the affected region are split into pre-alteration mutations 
$\alpha_{ij}$ and post-alteration mutations $\beta_{ij}$, yielding two equations:
%
\begin{align}
    \alpha_{ij} &= \mu \cdot w \cdot l_{CN\mathrm{tot},ij} \cdot c^\alpha_{ij} 
    \cdot (t_{CN,ij} - t_{\mathrm{pa}(i)}), \\[4pt]
    \beta_{ij} &= \mu \cdot w \cdot l_{CN\mathrm{tot},ij} \cdot c^\beta_{ij} 
    \cdot (t_i - t_{CN,ij}).
\end{align}
%
Solving the first equation for $t_{CN,ij}$ gives
%
\begin{equation}
    t_{CN,ij} = \frac{\alpha_{ij}}{\mu \cdot w \cdot l_{CN\mathrm{tot},ij} 
    \cdot c^\alpha_{ij}} + t_{\mathrm{pa}(i)},
\end{equation}
%
and substituting into the second equation gives
%
\begin{equation}
    t_i = \frac{\beta_{ij}}{\mu \cdot w \cdot l_{CN\mathrm{tot},ij} 
    \cdot c^\beta_{ij}} + t_{CN,ij}.
\end{equation}
%
Together, these equations form a linear system that can be solved sequentially 
along the tree, starting from the root and proceeding toward the leaves.

\subsection*{Example: Four Clones with Two Copy-Number Alterations}

Consider a tumour with four clones arranged in a linear topology,
%
\begin{equation}
    1 \longrightarrow 2 \longrightarrow 3 \longrightarrow 4,
\end{equation}
%
with $t_0 = 0$ and $T_f = 1$. Clone 2 carries a copy-number alteration event with 
timing $t_{CN,2}$, and clone 3 carries a copy-number alteration event with timing 
$t_{CN,3}$. The branches of clones 1 and 4 have no copy-number alterations. The 
full system of equations is as follows.

For the branch of clone 1, which has no copy-number alteration:
%
\begin{equation}
    m_1 = \mu \cdot w \cdot l_1 \cdot t_1 
    \quad \longrightarrow \quad 
    t_1 = \frac{m_1}{\mu \cdot w \cdot l_1}.
\end{equation}

For the branch of clone 2, which carries a copy-number alteration at $t_{CN,2}$:
%
\begin{align}
    \alpha_2 &= \mu \cdot w \cdot l_{CN,2} \cdot c^\alpha_2 
    \cdot (t_{CN,2} - t_1)
    \quad \longrightarrow \quad
    t_{CN,2} = \frac{\alpha_2}{\mu \cdot w \cdot l_{CN,2} \cdot c^\alpha_2} + t_1, 
    \\[4pt]
    \beta_2 &= \mu \cdot w \cdot l_{CN,2} \cdot c^\beta_2 
    \cdot (t_2 - t_{CN,2})
    \quad \longrightarrow \quad
    t_2 = \frac{\beta_2}{\mu \cdot w \cdot l_{CN,2} \cdot c^\beta_2} + t_{CN,2}.
\end{align}

For the branch of clone 3, which carries a copy-number alteration at $t_{CN,3}$:
%
\begin{align}
    \alpha_3 &= \mu \cdot w \cdot l_{CN,3} \cdot c^\alpha_3 
    \cdot (t_{CN,3} - t_2)
    \quad \longrightarrow \quad
    t_{CN,3} = \frac{\alpha_3}{\mu \cdot w \cdot l_{CN,3} \cdot c^\alpha_3} + t_2, 
    \\[4pt]
    \beta_3 &= \mu \cdot w \cdot l_{CN,3} \cdot c^\beta_3 
    \cdot (t_3 - t_{CN,3})
    \quad \longrightarrow \quad
    t_3 = \frac{\beta_3}{\mu \cdot w \cdot l_{CN,3} \cdot c^\beta_3} + t_{CN,3}.
\end{align}

For the branch of clone 4, which has no copy-number alteration:
%
\begin{equation}
    m_4 = \mu \cdot w \cdot l_4 \cdot (t_4 - t_3)
    \quad \longrightarrow \quad
    t_4 = \frac{m_4}{\mu \cdot w \cdot l_4} + t_3.
\end{equation}
%
The system can be solved sequentially: given the observed mutation counts 
$m_1$, $\alpha_2$, $\beta_2$, $\alpha_3$, $\beta_3$, $m_4$, the known mutation 
rate $\mu$, growth rate $w$, and genomic lengths $l$, all event times 
$t_1, t_{CN,2}, t_2, t_{CN,3}, t_3, t_4$ can be recovered exactly. The Poisson 
likelihoods in the main model represent the probabilistic counterpart of this 
deterministic system, where mutation counts are treated as noisy observations 
of their expected values.

\section{System of Equations for Unknown Growth Rates}
\label{app:RW}
\label{app:randomwalk}

When the growth rates are unknown, the deterministic system described in 
Appendix~\ref{app:KNWNO W} is no longer solvable in closed form, since each 
equation now contains two unknowns: the event time and the clone-specific growth 
rate $w_i$. Specifically, for a branch without copy-number alterations,
%
\begin{equation}
    m_i = \mu \cdot w_i \cdot l_{\mathrm{tot},i} \cdot (t_i - t_{\mathrm{pa}(i)}),
\end{equation}
%
and only the product $w_i \cdot (t_i - t_{\mathrm{pa}(i)})$ is identified by 
the data. Similarly, for branches with copy-number alterations,
%
\begin{align}
    \alpha_{ij} &= \mu \cdot w_i \cdot l_{CN\mathrm{tot},ij}
    \cdot c^\alpha_{ij} \cdot (t_{CN,ij} - t_{\mathrm{pa}(i)}), \\[4pt]
    \beta_{ij} &= \mu \cdot w_i \cdot l_{CN\mathrm{tot},ij} \cdot c^\beta_{ij} 
    \cdot (t_i - t_{CN,ij}),
\end{align}
%
where the clone growth rate $w_i$ is now itself unknown, so that each equation
mixes an event time with a growth rate. The system is therefore underdetermined,
and additional information is needed to separate timings from growth rates.

\subsection*{Example: Four Clones with Two Copy-Number Alterations and Unknown 
Growth Rates}

Consider the same four-clone linear topology as in Appendix~\ref{app:KNWNO W},
%
\begin{equation}
    1 \longrightarrow 2 \longrightarrow 3 \longrightarrow 4,
\end{equation}
%
but now with clone-specific growth rates $w_1, w_2, w_3, w_4$ treated as unknown. 
Both the pre- and post-alteration mutations on the branch of clone $i$ lie on the
stem lineage of that clone and therefore accumulate at its own rate $w_i$; the split
at $t_{CN,ij}$ changes only the number of allelic copies available for new mutations,
through $c^\alpha_{ij}$ and $c^\beta_{ij}$, not the rate. The full underdetermined
system is:
%
\begin{align}
    m_1 &= \mu \cdot w_1 \cdot l_1 \cdot t_1, \\[4pt]
    \alpha_2 &= \mu \cdot w_2 \cdot l_{CN,2} \cdot c^\alpha_2
    \cdot (t_{CN,2} - t_1), \\[4pt]
    \beta_2 &= \mu \cdot w_2 \cdot l_{CN,2} \cdot c^\beta_2 
    \cdot (t_2 - t_{CN,2}), \\[4pt]
    \alpha_3 &= \mu \cdot w_3 \cdot l_{CN,3} \cdot c^\alpha_3
    \cdot (t_{CN,3} - t_2), \\[4pt]
    \beta_3 &= \mu \cdot w_3 \cdot l_{CN,3} \cdot c^\beta_3 
    \cdot (t_3 - t_{CN,3}), \\[4pt]
    m_4 &= \mu \cdot w_4 \cdot l_4 \cdot (t_4 - t_3).
\end{align}
%
This system has $6$ equations in $10$ unknowns: the six event times
$t_1, t_{CN,2}, t_2, t_{CN,3}, t_3, t_4$ and the four growth rates
$w_1, w_2, w_3, w_4$. It is therefore underdetermined by $4$ degrees of freedom,
and no unique solution exists from the mutation data alone.

\subsection*{The Abundance Constraint and the Residual Scale}

The count above uses only the mutation data. The model, however, also observes the
clone abundances, and these add further constraints. Suppose first that the
\emph{absolute} size of each clade were observed. The terminal mass of clone $i$
would then satisfy
%
\begin{equation}
    N_i = e^{\,w_i\,(T_f - t_i)}
    \quad\Longrightarrow\quad
    \ln N_i = w_i\,(T_f - t_i),
\end{equation}
%
one equation per clone. Combined with the mutation equation
$m_i = \mu \cdot w_i \cdot l_i \cdot (t_i - t_{\mathrm{pa}(i)})$, the two time
intervals telescope,
%
\begin{equation}
    \frac{m_i}{\mu\,l_i} + \ln N_i
    = w_i\,(t_i - t_{\mathrm{pa}(i)}) + w_i\,(T_f - t_i)
    = w_i\,(T_f - t_{\mathrm{pa}(i)}),
\end{equation}
%
which solves for the growth rate and then for the event time,
%
\begin{equation}
    w_i = \frac{m_i/(\mu\,l_i) + \ln N_i}{T_f - t_{\mathrm{pa}(i)}},
    \qquad
    t_i = t_{\mathrm{pa}(i)} + \frac{m_i}{\mu\,l_i\,w_i}.
\end{equation}
%
With absolute clade sizes the system would therefore be \emph{exactly identified}:
each clone contributes one mutation equation and one abundance equation, matching its
two unknowns $w_i$ and $t_i$ (on a copy-number branch the extra equation $\beta_{ij}$
matches the extra unknown $t_{CN,ij}$).

The single-cell data, however, do not provide absolute sizes. Sequencing reads a
number of cells $N_{\mathrm{tot}}$ fixed by the experiment, not the true tumour size,
so what is observed is
%
\begin{equation}
    N_i = c \cdot e^{\,w_i\,(T_f - t_i)},
    \qquad
    c = \frac{N_{\mathrm{tot}}}{\sum_k M_k},
\end{equation}
%
with an unknown common sampling factor $c$; equivalently, only the proportions
$p_i = M_i / \sum_k M_k$ are measured. Taking logarithms,
%
\begin{equation}
    \ln N_i = w_i\,(T_f - t_i) + \ln c,
\end{equation}
%
so the abundances determine $w_i\,(T_f - t_i)$ only up to the common constant
$\ln c$: they fix the $K-1$ \emph{ratios} of clade sizes but not their overall scale.
In the four-clone example these three independent ratios cut the under-determination
from four down to one; in general the $K$ mutation equations and the $K-1$ abundance
ratios pin the $2K$ growth-and-time unknowns up to a single residual degree of freedom
(each copy-number branch adds one timing unknown and one equation, leaving the balance
unchanged). That one remaining direction is precisely the global growth scale $c$,
which a single snapshot cannot resolve because the composition of one sample is
invariant to it. It is this residual scale, rather than four free directions, that the
hierarchical prior must fix, as discussed next; recovering it from the data instead
would require an absolute measure of tumour burden, ideally collected at more than one
time.

\subsection*{Hierarchical Prior as Regularisation}

The additional information needed to regularise the system is provided by the 
hierarchical prior on the growth rates described in Section~\ref{sec: model rw}. 
The random walk on the log scale assigns
%
\begin{align}
    \log w_1 &= \log \bar{w}, \\[4pt]
    \log w_2 &= \log w_1 + \sigma_{\log w} \cdot z_2, \\[4pt]
    \log w_3 &= \log w_2 + \sigma_{\log w} \cdot z_3, \\[4pt]
    \log w_4 &= \log w_3 + \sigma_{\log w} \cdot z_4,
\end{align}
%
with $z_2, z_3, z_4 \overset{\mathrm{iid}}{\sim} \mathcal{N}(0,1)$ and 
$\log \bar{w} \sim \mathcal{N}(\log \bar{w}_{\mathrm{prior}}, \sigma_{\bar{w}})$. 
This introduces $4$ additional parameters ($\bar{w}, z_2, z_3, z_4$) governed 
by prior distributions that constrain their values, effectively regularising the 
posterior. The prior encodes the biological assumption that related clones tend 
to have similar growth rates: the further two clones are apart in the tree, the 
more their rates are allowed to differ, since deviations accumulate along the 
branches.

In the limit $\sigma_{\log w} \to 0$, all growth rates collapse to $\bar{w}$ and
the system reduces to the identifiable case of Appendix~\ref{app:KNWNO W}, where
$w = \bar{w}$ is known. For finite $\sigma_{\log w}$, the prior allows the clone
rates to differ while penalising large departures.

It is worth being precise about what this regularisation does and does not achieve.
It makes the posterior \emph{proper} and lets the sampler explore plausible
combinations of timings and growth rates, but it does not make the growth rates
\emph{identified by the data}. Along the single residual direction identified in the
previous subsection, the global growth scale $c$, both the mutation likelihood and the
abundance proportions are unchanged, so there the posterior simply follows the prior.
The residual uncertainty this leaves on the
per-clone rates is exactly the single-snapshot floor observed in the simulation
study, and it cannot be removed by sharpening the prior: a tighter prior would narrow
the posterior only by asserting information the data do not provide, at the cost of
calibration. What \emph{is} learned from the data is the degree of pooling toward a
common rate, through the posterior of $\sigma_{\log w}$; the overall growth scale, by
contrast, is set by the prior rather than by the single snapshot.